\clearpage
\subsection{Perceived AI Support — AI Questionnaire}
\label{sec:eval-ai-support}

The SUS (Section~\ref{sec:eval-sus}) and R-TLX
(Section~\ref{sec:eval-rtlx}) characterize the interface as it was
experienced; they say nothing directly about whether participants
felt the AI itself had been useful. The Perceived AI Support
questionnaire targets that gap. After each interaction mode,
participants completed a six-item questionnaire on a five-point
Likert scale ($1$ = strong disagreement, $5$ = strong agreement):
helpfulness (item~1), efficiency (item~2), transparency (item~3),
trust (item~4), legal-basis confidence (item~5), and re-use intent
(item~6). Item~5 is the primary self-report confidence measure
referenced in Chapter~\ref{ch:methodology}. The full per-item
distribution is in Appendix~\ref{app:items-ai}.

\begin{table}[htbp]
\centering
\caption{Perceived AI Support, mean response per item by mode
($n = 69$ per mode), five-point Likert scale (higher is better).
The Aggregated column is the unweighted mean of the three per-mode
means.}
\label{tab:ai-per-item}
\small
\begin{tabularx}{\linewidth}{@{}c X r r r r@{}}
\toprule
\# & Item & Aggregated & Form & Ask & Chat \\
\midrule
1 & The AI was helpful in creating this document. & 4.23 & 4.23 & 4.32 & 4.13 \\
2 & The AI helped me to complete the document faster. & 4.22 & 4.43 & 4.25 & 3.99 \\
3 & I understood why the AI suggested specific GDPR fields. & 3.81 & 3.51 & 4.16 & 3.75 \\
4 & I trust that the AI made correct suggestions. & 3.63 & 3.41 & 3.80 & 3.68 \\
5 & I am confident that the AI identified the correct legal basis. & 3.71 & 3.58 & 3.74 & 3.81 \\
6 & I would use this mode again. & 3.65 & 3.99 & 3.55 & 3.42 \\
\bottomrule
\end{tabularx}
\end{table}

\begin{figure}[htbp]
\centering
\includegraphics[width=0.45\linewidth]{ai_item_1_helpful.png}\hfill
\includegraphics[width=0.45\linewidth]{ai_item_2_efficiency.png}\\[1ex]
\includegraphics[width=0.45\linewidth]{ai_item_3_transparency.png}\hfill
\includegraphics[width=0.45\linewidth]{ai_item_4_trust.png}\\[1ex]
\includegraphics[width=0.45\linewidth]{ai_item_5_legal.png}\hfill
\includegraphics[width=0.45\linewidth]{ai_item_6_reuse.png}
\caption{Box plots of the six Perceived AI Support items by mode
($n = 69$ per mode), one plot per item in administered order.}
\label{fig:ai-subquestions}
\end{figure}

Asked simply whether the AI was helpful in creating the document,
participants rated all three modes similarly: Ask drew the highest
rating ($4.32$), Form just behind ($4.23$), and Chat lowest
($4.13$), a spread of only $0.19$ points. The AI is perceived as
broadly helpful across all three modes; the differences emerge
only when participants are asked more specifically what the AI was
helpful \emph{for}.

On whether the AI helped complete the document faster, Form drew
the strongest endorsement ($4.43$), Ask second ($4.25$), and Chat
lowest ($3.99$). The perceived speed-up decreases monotonically as
LLM involvement grows, dropping by roughly $0.45$ points from Form
to Chat. Participants felt the AI helped them most when it stayed
in the background.

Participants felt they understood the AI's suggestions best in Ask
($4.16$), well ahead of Chat ($3.75$) and Form ($3.51$). The Ask
lead of roughly $0.65$ points is the widest reversal of the
Form-leading pattern across the six items, and is unsurprising:
Ask is the only mode that explicitly explains each suggestion (see
Section~\ref{sec:ropagen}).

Trust in the AI's outputs follows the same direction. Asked
whether they trusted that the AI made correct suggestions,
participants rated Ask highest ($3.80$), Chat second ($3.68$), and
Form lowest ($3.41$). The LLM-driven modes earn more trust than
Form, by roughly $0.27$ to $0.39$ points.

On confidence that the AI identified the correct legal basis (the
primary confidence measure, item 5), the ordering reverses
completely: Chat drew the highest rating ($3.81$), Ask second
($3.74$), and Form lowest ($3.58$). Confidence rises monotonically
with LLM involvement, the only item where Chat leads outright.
Participants feel most sure of the LLM's legal judgment in the
mode where they see the LLM least directly.

Asked which mode they would actually use again, participants
snapped back to Form ($3.99$) over Ask ($3.55$) and Chat ($3.42$),
a $0.57$-point lead. This is the widest negative gap of the six
items and the only one to reach statistical significance
($\chi^2(2) = 11.12$, $p = 0.004$; Form vs Chat $r = +0.543$,
large). Even when participants felt better about the LLM's content
judgment, they still wanted to come back to Form.

The composite Friedman test on the per-mode aggregate is not
significant ($\chi^2(2) = 1.05$, $p = 0.592$, $W = 0.008$), so
perceived AI support taken as a whole does not differ across
modes. The primary legal-basis confidence item (item 5) is also
non-significant ($\chi^2(2) = 3.76$, $p = 0.152$). Full statistics
in Appendix~\ref{app:inferential}.

Perceived AI support splits the user-study story in two:
participants credit the LLM-driven modes when asked about content
judgment (transparency, trust, legal-basis confidence), but credit
Form when asked about productivity and re-use intent. The single
item that reaches statistical significance after Bonferroni
correction (re-use intent, item 6) is the one where the Form lead
is sharpest: participants are willing to praise the LLM modes'
content judgment, but still want to come back to Form. The
behavioral counterpart, which mode they actually choose when
forced to rank, is taken up next in
Section~\ref{sec:eval-freeform}.
